% The pipeline drawing itself. \input by figures/fig_pipeline.tex for the
% typeset build, and by src/make_pipeline_png.py, which renders this same
% picture to figures/fig_pipeline.png for the Word build. Edit only here, so
% the two builds cannot drift apart.
%
% Flow: message -> model -> verdict; only the malicious verdicts become the
% attack set; the perturbation engine edits them under the shared budget and
% re-queries the model (the loop back to row 1); the stored per-sample records
% are what the transfer and defence evaluation reads.
\begin{tikzpicture}[
  box/.style   = {draw=black!65, line width=0.6pt, fill=black!4,
                  rounded corners=2pt, align=center, font=\footnotesize,
                  inner sep=2.5pt, text width=1.95cm, minimum height=1.05cm},
  arrow/.style = {-{Latex[length=1.5mm, width=1.1mm]}, draw=black!75,
                  line width=0.6pt},
  note/.style  = {font=\scriptsize, text=black!70, align=center},
]
  \node[box] (msg) {Input\\message};
  \node[box, right=8mm of msg] (model) {Language model\\(triage)};
  \node[box, right=8mm of model] (verdict) {Verdict\\(malicious\\ or benign)};

  \node[box, below=11mm of verdict] (set) {Attack set\\(caught messages)};
  \node[box, left=8mm of set] (engine) {Perturbation engine\\(15\% budget)};
  \node[box, left=8mm of engine] (eval) {Transfer \& defence\\evaluation};

  \draw[arrow] (msg) -- (model);
  \draw[arrow] (model) -- (verdict);
  \draw[arrow] (verdict) -- node[note, right=1.2mm, text width=1.05cm]
      {malicious\\only} (set);
  \draw[arrow] (set) -- node[note, above=1.4mm, text width=1.1cm]
      {caught\\text} (engine);
  % The label sits below the box band: between the boxes there is only the 8 mm
  % gap, which a word does not fit into without touching either border.
  \draw[arrow] (engine) -- node[note, below=5.7mm, text width=1.0cm]
      {records} (eval);
  \draw[arrow] (engine) -- node[note, left=1.2mm, text width=1.15cm]
      {edited text,\\re-query} (model);
\end{tikzpicture}
